\exerciseeditionsection{Solutions to Wald's Problems}
\ExerciseEditorialNote

\begin{enumerate}[leftmargin=2.8em,itemsep=1.2em]
\WaldSolution{1}
\begin{enumerate}[label=\alph*),leftmargin=2em]
\item Write the three-metric as
\(\dd s^2=A(r)\dd r^2+r^2\dd\Omega^2\).  Matching it to
\(H(\bar r)(\dd\bar r^2+\bar r^2\dd\Omega^2)\) requires
\[
 H=\frac{r^2}{\bar r^2},
 \qquad
 \frac{\dd\bar r}{\bar r}=\frac{\sqrt{A(r)}}{r}\,\dd r.
\]
The second equation has a local solution wherever \(\nabla_ar\ne0\), proving
the conformally flat form.

\item For Schwarzschild,
\(A=(1-2M/r)^{-1}\).  Integration and the asymptotic choice
\(\bar r\sim r\) give
\[
 r=\bar r\left(1+\frac{M}{2\bar r}\right)^2.
\]
Substitution into the temporal and spatial terms yields
\[
 \dd s^2=-\frac{(1-M/2\bar r)^2}{(1+M/2\bar r)^2}\dd t^2
 +\left(1+\frac{M}{2\bar r}\right)^4
 (\dd\bar r^2+\bar r^2\dd\Omega^2).
\]
\end{enumerate}

\WaldSolution{2}
For
\(\dd s^2=-f\dd t^2+h\dd r^2+r^2\dd\Omega^2\), the independent nonzero
coordinate components are
\[
\begin{aligned}
R_{tt}&=\frac{f}{h}\left(
 \frac{f''}{2f}-\frac{f'^2}{4f^2}-\frac{f'h'}{4fh}
 +\frac{f'}{fr}\right),\\
R_{rr}&=-\frac{f''}{2f}+\frac{f'^2}{4f^2}
 +\frac{f'h'}{4fh}+\frac{h'}{hr},\\
R_{\theta\theta}&=1-\frac1h+\frac{rh'}{2h^2}
 -\frac{rf'}{2fh},\\
R_{\phi\phi}&=\sin^2\theta\,R_{\theta\theta}.
\end{aligned}
\]
All mixed components vanish.  The coordinate calculation requires the full
Christoffel list and several cancellations; the tetrad calculation packages
the same work into a few connection one-forms.

\WaldSolution{3}
\begin{enumerate}[label=\alph*),leftmargin=2em]
\item Static spherical symmetry leaves one invariant two-form on the
\((t,r)\) plane and the area form on each orbit sphere.  Their coefficients
depend only on \(r\), giving exactly the stated \(A(r)\) and \(B(r)\) form.

\item With \(B=0\), the homogeneous equation is automatic.  The sourced-free
equation reduces to
\(\partial_r(r^2A)=0\), so \(A=C/r^2\).  With Wald's orientation and charge
convention, Gauss's law fixes \(C=-q\).

\item In the orthonormal frame the electric stress tensor is
\[
 T_{\hat\mu\hat\nu}=\frac{q^2}{8\pi r^4}
 \operatorname{diag}(1,-1,1,1).
\]
Write \(h=(1-2m(r)/r)^{-1}\).  The \(tt\) equation gives
\(m'=q^2/(2r^2)\), hence
\(m=M-q^2/(2r)\).  The difference of the \(tt\) and \(rr\) equations gives
\((fh)'=0\); asymptotic flatness fixes \(fh=1\).  Therefore
\[
 f=h^{-1}=1-\frac{2M}{r}+\frac{q^2}{r^2},
\]
which is the Reissner--Nordstr\"om metric.
\end{enumerate}

\WaldSolution{4}
\begin{enumerate}[label=\alph*),leftmargin=2em]
\item A stationary observer has \(u^a=\xi^a/V\), and
\(\xi^a\nabla_aV=0\).  The Killing equation gives
\(\xi^a\nabla_a\xi_b=-\tfrac12\nabla_b(\xi^2)=V\nabla_bV\).
Thus
\[
 a_b=u^a\nabla_au_b=\nabla_b\ln V.
\]

\item A stationary particle has energy at infinity \(E=mV\).  A local
outward displacement \(\dd\ell\) changes it by
\[
 \dd E=m\,\dd V=m|\nabla V|\dd\ell=VF\,\dd\ell.
\]
Quasistatic energy conservation equates this with the work done at the remote
end of the massless string.  Hence \(F_\infty=VF\).
\end{enumerate}

\WaldSolution{5}
For an equatorial Schwarzschild null geodesic, conservation of energy and
angular momentum gives
\[
 \frac{\dd t}{\dd r}
 =\frac1c\left(1-\frac{2GM}{c^2r}\right)^{-1}
 \left[1-\frac{b^2}{r^2}
 \left(1-\frac{2GM}{c^2r}\right)\right]^{-1/2}.
\]
At closest approach \(r=R_0\),
\[
 b^2=\frac{R_0^2}{1-2GM/(c^2R_0)}.
\]
Insert this relation, expand the integrand through first order in \(GM/c^2\),
and integrate from \(R_0\) to \(R_E\) and from \(R_0\) to \(R_P\).  Doubling
for the return trip gives
\[
\begin{aligned}
\Delta t=\frac{2}{c}\biggl\{&\sqrt{R_E^2-R_0^2}+\sqrt{R_P^2-R_0^2}\\
&+\frac{2GM}{c^2}\biggl[
2\ln\frac{R_E+\sqrt{R_E^2-R_0^2}}{R_0}
+2\ln\frac{R_P+\sqrt{R_P^2-R_0^2}}{R_0}\\
&\hspace{3.2em}+\sqrt{\frac{R_E-R_0}{R_E+R_0}}
+\sqrt{\frac{R_P-R_0}{R_P+R_0}}
\biggr]\biggr\},
\end{aligned}
\]
which is Eq.~\eqref{eq:6.3.45}.

\WaldSolution{6}
In region II put \(A=2M/r-1>0\).  For any timelike world line,
normalization gives
\[
 -1=A\dot t^2-A^{-1}\dot r^2
 +r^2(\dot\theta^2+\sin^2\theta\dot\phi^2).
\]
Therefore \(A^{-1}\dot r^2\ge1\), or
\[
 |\dot r|\ge\sqrt{2M/r-1}.
\]
Future direction in region II selects decreasing \(r\).  The largest possible
proper time is consequently
\[
 \tau_{\max}=\int_0^{2M}\frac{\dd r}{\sqrt{2M/r-1}}.
\]
With \(r=2M\sin^2\chi\), this becomes
\(4M\int_0^{\pi/2}\sin^2\chi\,\dd\chi=\pi M\).
Equality requires \(\dot t=\dot\theta=\dot\phi=0\), the limiting radial
geodesic entering from \(r=2M\) with conserved energy \(E\to0\).
\end{enumerate}
